\documentclass[11pt]{article}
\usepackage[margin=1in]{geometry}
\usepackage{booktabs}
\usepackage{amsmath,amssymb}
\usepackage{graphicx}
\usepackage{hyperref}

% Paper-facing assets are copied through tools/export_manuscript_assets.py,
% which records their source paths and hashes.
\graphicspath{
  {../generated/figures/adaptive_h/}
}

\title{Final Scale-Adaptive CKME-DCP Experiment Report}
\author{}
\date{July 23, 2026}

\begin{document}
\maketitle

\paragraph{Status.}
This report records the completed 50-macroreplication final benchmark. The run
uses iid target-law calibration with one response per input, excludes the
diagnostic \(S^0\) score, and compares fixed, oracle, and sample-SD plus
Nadaraya--Watson bandwidths on paired data. The final QA status is PASS with one
reporting-layer warning described below.

\section{Research question}

The standard CKME conditional CDF estimator uses a single CV-tuned smooth-indicator
bandwidth. Under heterogeneous conditional response scales, the effective
bandwidth is \(h/s(x)\), so a fixed \(h\) can oversmooth low-scale regions and
undersmooth high-scale regions. We ask two separate questions:
\begin{enumerate}
  \item Does replacing the scalar bandwidth with the oracle rule
\[
  h(x)=c\,s(x)
\]
  reduce raw-score heterogeneity or improve interval performance?
  \item Does the implementable rule \(h(x)=c\widehat s(x)\), with
  \(\widehat s\) estimated from replicated Stage-1 data, retain that behavior?
\end{enumerate}

\section{Locked setup}

The three DGPs are an exact stationary M/M/1 sojourn time,
\[
  Y\mid\rho\sim\operatorname{Exponential}(\text{rate}=1-\rho),
  \qquad \rho\in[0.1,0.9],
\]
and two raised-floor location-scale models with mean
\(e^{x/10}\sin x\), scale \(s(x)=0.10+0.20(x-\pi)^2\), and either Gaussian
noise or variance-normalized \(t_3/\sqrt{3}\) noise.

Stage~1 uses \(r_0=10\) and
\(B=n_0r_0\in\{100,250,500,1000\}\). Each macroreplication uses 1,000 iid
calibration pairs and 1,000 iid test pairs from the same uniform target input
law, with one response per input. We use 50 macroreplications, \(c=1\), and a
frozen Silverman factor of one for the Nadaraya--Watson input bandwidth. All
three arms share Stage-1, calibration, test, and seed streams.

Raw point scores define coverage. A running-max projected CDF and generalized
inverse define only the reported interval, projected-interval coverage, width,
and interval score. The fixed scalar bandwidth is selected once on a separate
site-grouped pretraining sample. The plug-in multiplier and Nadaraya--Watson
factor are frozen at one rather than comparably tuned, so this is a
frozen-default mechanism benchmark.

\section{Main results at \texorpdfstring{\(B=1000\)}{B=1000}}

\begin{table}[p]
\centering
\caption{Final adaptive-bandwidth benchmark at $B=1000$. Parentheses report Monte Carlo standard errors.}
\label{tab:final-adaptive-h}
\textbf{(a) Coverage and score-set/interval agreement}\par\smallskip
\small
\begin{tabular}{@{}llcccc@{}}
\toprule
DGP & Method & Raw cov. (MCSE) & Interval cov. (MCSE) & Disagree. & Worst gap \\
\midrule
M/M/1 & Fixed & 0.899 (0.002) & 0.855 (0.002) & 0.049 & 0.072 \\
M/M/1 & Oracle & 0.899 (0.002) & 0.902 (0.002) & 0.004 & 0.062 \\
M/M/1 & Plug-in & 0.898 (0.002) & 0.903 (0.002) & 0.004 & 0.068 \\
Gaussian & Fixed & 0.901 (0.002) & 0.900 (0.002) & 0.003 & 0.112 \\
Gaussian & Oracle & 0.900 (0.002) & 0.900 (0.002) & 0.005 & 0.084 \\
Gaussian & Plug-in & 0.902 (0.002) & 0.902 (0.002) & 0.003 & 0.127 \\
Student-$t_3$ & Fixed & 0.902 (0.002) & 0.902 (0.002) & 0.007 & 0.112 \\
Student-$t_3$ & Oracle & 0.900 (0.002) & 0.900 (0.002) & 0.009 & 0.088 \\
Student-$t_3$ & Plug-in & 0.901 (0.002) & 0.901 (0.002) & 0.006 & 0.117 \\
\bottomrule
\end{tabular}
\par\medskip
\textbf{(b) Projected-interval efficiency}\par\smallskip
\begin{tabular}{@{}llcc@{}}
\toprule
DGP & Method & Width (MCSE) & Interval score (MCSE) \\
\midrule
M/M/1 & Fixed & 7.904 (0.057) & 10.983 (0.097) \\
M/M/1 & Oracle & 7.221 (0.063) & 12.463 (0.133) \\
M/M/1 & Plug-in & 7.020 (0.055) & 12.568 (0.147) \\
Gaussian & Fixed & 2.420 (0.014) & 3.334 (0.018) \\
Gaussian & Oracle & 2.454 (0.014) & 3.198 (0.018) \\
Gaussian & Plug-in & 2.296 (0.013) & 3.372 (0.024) \\
Student-$t_3$ & Fixed & 2.051 (0.020) & 3.576 (0.038) \\
Student-$t_3$ & Oracle & 1.923 (0.016) & 3.373 (0.037) \\
Student-$t_3$ & Plug-in & 1.842 (0.015) & 3.486 (0.041) \\
\bottomrule
\end{tabular}
\par\medskip
\begin{minipage}{0.98\linewidth}
\footnotesize\textit{Notes.} Raw-score coverage is guarantee-bearing. Interval coverage, width, and interval score use the monotone-projected generalized-inverse interval. Disagreement is the fraction of test points for which raw score-set membership and projected-interval membership differ. Worst group gaps use raw-score coverage in ten equal-count input bins. Gaussian and Student-$t_3$ denote the two raised-floor DGPs; plug-in denotes sample-SD plus Nadaraya--Watson smoothing.
\end{minipage}
\end{table}


All nine raw-score marginal coverages are between 0.898 and 0.902. The oracle
mechanism is strong on the two raised-floor DGPs: maximum pairwise score KS
falls from 0.531 to 0.296 under Gaussian noise and from 0.718 to 0.437 under
Student-\(t_3\) noise. Their oracle worst-group gaps also fall from 0.112 to
0.084 and 0.088.

The raw-score guarantee does not automatically transfer to the projected
interval. For fixed-bandwidth M/M/1, raw coverage is 0.899 but projected
interval coverage is 0.855, with 0.049 score-set/interval disagreement. The
other eight \(B=1000\) combinations have projected coverage from 0.900 to
0.903 and disagreement from 0.003 to 0.009. Width and interval score must
therefore be interpreted as reporting-layer metrics.

The plug-in result is mixed. It produces narrower intervals than fixed
bandwidth on all three DGPs, but improves interval score only on
Student-\(t_3\): 3.486 versus 3.576. Its interval score is 3.372 versus 3.334
for Gaussian noise and 12.568 versus 10.983 for M/M/1. Its worst-group gaps do
not uniformly retain the oracle improvement.

\section{Scale and score diagnostics}

The plug-in mean absolute relative scale error decreases from budget 100 to
1,000: 0.182 to 0.075 for M/M/1, 0.920 to 0.353 for Gaussian noise, and 0.817
to 0.308 for Student-\(t_3\). Thus the estimator increasingly, but only
partially, tracks the scale as Stage-1 information grows. This raw error
includes the fixed-\(r_0\) global-multiplier difference and is not a direct
test of uniform scale-pattern consistency.

\begin{figure}[ht!]
\centering
\includegraphics[width=\linewidth]{final_effective_bandwidth_ratio.pdf}
\caption{Effective bandwidth ratio at \(B=1000\).}
\end{figure}

On the two raised-floor models, the plug-in score KS lies between fixed and
oracle at \(B=1000\): 0.468 for Gaussian noise and 0.539 for
Student-\(t_3\). M/M/1 is the boundary case: plug-in KS is 0.286, compared
with 0.224 for fixed and 0.231 for oracle. Because this statistic is the
maximum of 45 finite-sample pairwise KS distances, it is a relative mechanism
diagnostic, not a null-calibrated test of exact homogeneity.

\begin{figure}[ht!]
\centering
\includegraphics[width=\linewidth]{final_raw_score_homogeneity.pdf}
\caption{Maximum pairwise raw-score KS distance across ten equal-count input
bins. Error bars are Monte Carlo standard errors.}
\end{figure}

\section{QA and safe interpretation}

The audit verifies 600 jobs, 1,800 per-point files, complete paired arms,
identical test data, finite outputs, complete ten-bin score diagnostics, and
hash-bound publication assets. The one warning is confined to reporting
interval inversion: two of 1,800 arm files exceed a 2\% grid-boundary rate,
both fixed-bandwidth \(B=100\) runs. The mean rate is \(2.14\times10^{-4}\),
and no \(B=1000\) arm has a boundary hit. Raw-score coverage is unaffected.

The safe conclusion is:
\begin{quote}
Oracle scale normalization reduces score heterogeneity in the raised-floor
Gaussian and heavy-tailed models. The sample-SD plus Nadaraya--Watson plug-in
increasingly tracks the scale and moves the score diagnostic toward the oracle
in those models, but its interval-score and groupwise gains are DGP-dependent.
\end{quote}

The results do not support a claim that the plug-in uniformly dominates fixed
bandwidth. The main next questions are selection of the global multiplier,
refinement of the scale smoother, projection-aware interval construction, a
null reference for the KS diagnostic, and the smallest external baseline set
needed for a journal comparison.

\end{document}
